\chapter{Software Development Process}

This chapter describes the software development process adopted for the AI-Driven Smart Home
Energy Optimizer (SHEO) project. It grounds that description in the theoretical framework
taught in the course: the software development life cycle (SDLC), the classical process
models, the agile philosophy, and the ISO/IEC~9126 quality model. Each section moves from
the general course concept to its concrete instantiation in SHEO.

% ────────────────────────────────────────────────────────────────────────────────
\section{The Software Development Life Cycle}
\label{sec:sdlc}

A software development life cycle (SDLC) partitions the construction of a software system
into a sequence of well-defined, measurable phases, each with assigned responsibilities and
verifiable outputs. The course identifies two complementary views of the SDLC.

\textbf{Phase-level view.} Five core phases characterise most plan-driven processes:
\textit{Requirements Analysis} (understand and specify what the system must do),
\textit{Design} (determine how it will be structured),
\textit{Development / Coding} (realise the design in executable form),
\textit{Testing} (verify correctness and validate fitness for use), and
\textit{Maintenance} (correct, adapt, and extend the deployed system).

\textbf{Generic project-step view.} At the project level these phases are bounded by two
additional activities: a \textit{Feasibility Study} at the outset (to establish that the
problem is worth solving and is tractable within the given constraints) and
\textit{Acceptance \& Release / Operation} at the close. The full generic sequence is
therefore:

\begin{center}
\textit{Feasibility Study}
$\rightarrow$ \textit{Requirements}
$\rightarrow$ \textit{UI Design}
$\rightarrow$ \textit{System Design}
$\rightarrow$ \textit{Program Development}
$\rightarrow$ \textit{Acceptance \& Release}
$\rightarrow$ \textit{Operation \& Maintenance}
\end{center}

The SHEO deliverable set mirrors this sequence directly. A feasibility study assessed
technical and economic viability and fixed the scope decision to simulate hardware rather
than depend on physical IoT devices. The Software Requirements Specification formalised
the \textit{what}; the Software Design Document (SDD) recorded the \textit{how}; the
backend codebase and React frontend constituted the \textit{Development} phase; the
automated test suite and test plan covered \textit{Testing}; and a configuration management
plan addressed ongoing \textit{Maintenance} concerns. Crucially, the SDLC is not merely a
checklist: quality control operated continuously across all phases rather than being
deferred to a final testing stage, as the ISO/IEC~9126 standard prescribes.

% ────────────────────────────────────────────────────────────────────────────────
\section{Comparison of Process Models}
\label{sec:proc-models}

The course presents six classical process models. Table~\ref{tab:process-models} summarises
the core idea and the conditions under which each model is most appropriate.

\begin{longtable}{@{}p{3.2cm} p{5.5cm} p{5.5cm}@{}}
\caption{Comparison of software process models}
\label{tab:process-models} \\
\toprule
\textbf{Model} & \textbf{Core idea} & \textbf{Best suited when} \\
\midrule
\endfirsthead
\multicolumn{3}{c}{\tablename~\ref{tab:process-models} (continued)} \\
\toprule
\textbf{Model} & \textbf{Core idea} & \textbf{Best suited when} \\
\midrule
\endhead
\bottomrule
\endfoot
\bottomrule
\endlastfoot
Waterfall &
  Sequential phases; each phase must be completed and signed off before the next begins.
  No iteration once a phase closes. &
  Requirements are fully understood and stable from the outset; domain is mature and
  well-documented; regulatory or safety standards mandate formal sign-off at each gate. \\
\addlinespace
Modified Waterfall &
  Adds limited feedback arrows between adjacent phases so that a discovered defect may be
  corrected in the preceding phase without restarting the entire process. &
  Requirements are largely stable but minor clarifications are anticipated; a more
  disciplined team than pure Waterfall, but still primarily sequential. \\
\addlinespace
Prototyping &
  A quick, often throw-away prototype is built to elicit or validate requirements before
  full development begins. Two variants: \textit{throw-away} (discard after learning) and
  \textit{evolutionary} (grow into the product). &
  Requirements are poorly understood or highly volatile; user interface expectations are
  hard to specify without tangible feedback; exploratory research context. \\
\addlinespace
Incremental &
  The system is partitioned into increments, each of which delivers usable functionality.
  The core system is built first; subsequent increments add capability on a working base. &
  A working subset of the system can be deployed and used early; requirements are
  moderately understood; a small to medium team; value is best delivered in slices. \\
\addlinespace
RAD &
  Rapid Application Development uses intensive, time-boxed cycles (typically 60--90~days)
  with heavy reuse of components and automated tooling to compress the schedule. &
  Time to market is critical; experienced team with strong tool knowledge; requirements can
  be scoped tightly enough for short cycles; reusable components are available. \\
\addlinespace
Spiral &
  Each iteration traverses four quadrants: planning, risk analysis, engineering, and
  evaluation. Risk is the primary driver for deciding whether to proceed, prototype, or
  abandon. &
  Large, complex, high-risk systems; risk identification and mitigation are primary
  concerns; budget exists for repeated analysis cycles; safety-critical or
  mission-critical domains. \\
\end{longtable}

% ────────────────────────────────────────────────────────────────────────────────
\section{The Process Model Adopted for SHEO}
\label{sec:adopted}

\subsection{Model selection}

SHEO was developed using an \textbf{incremental, agile-style process}. This is closest
to the Incremental model in the course taxonomy, augmented by agile working practices
described in Section~\ref{sec:agile}. The selection was justified by the following
conditions, cross-referenced against the course's model-selection criteria.

\begin{enumerate}[leftmargin=*, label=(\arabic*)]
  \item \textbf{Moderately understood requirements.} The functional scope---energy
        monitoring, capability-driven device control, rule-based automation, and
        AI-generated recommendations---was clear at a high level from the outset, but many
        interface details and edge-case safety constraints emerged during development. This
        ruled out pure Waterfall (which demands fully stable requirements) and favoured an
        iterative approach in which each increment surfaced and resolved ambiguities in a
        bounded scope.

  \item \textbf{Small team, short horizon.} A student team working over a single semester
        cannot sustain the overhead of formal Spiral risk-quadrant analysis, nor the
        intensive reuse infrastructure that RAD presupposes. An incremental model with
        lightweight process keeps coordination costs proportionate to team size.

  \item \textbf{Value delivered in working slices.} An early runnable system---even with
        only authentication, device registration, and a simulated sensor feed---provided a
        concrete foundation for demonstrating correctness and gathering feedback. Each
        subsequent increment extended a tested, running product rather than an accumulation
        of paper designs.

  \item \textbf{Reproducible demonstration requirement.} The project required a seeded demo
        that could be reliably executed at any point. This is best served by always
        maintaining a runnable product, which the incremental approach enforces naturally.
\end{enumerate}

\subsection{Incremental construction of SHEO}

Development proceeded through five broad increments, each leaving the system in a passing,
demonstrable state before the next began.

\begin{enumerate}[leftmargin=*, label=\textbf{Increment \arabic{enumi}.}]
  \item \textbf{Runnable spine.} Authentication (token-based), the user-role model
        (Administrator / Resident / Developer), device registration, and the database
        schema. The test suite was initialised; the seeded demo was executable from this
        point.

  \item \textbf{Capability schema and simulator.} The declarative capability schema---which
        describes every device's controllable parameters, constraints, and UI contract---was
        introduced, together with the background simulator that drives synthetic sensor
        readings. All existing tests continued to pass.

  \item \textbf{Energy monitoring and safety rules.} Real-time energy tracking endpoints,
        the safety-rule enforcement layer (NFR-SAF: rate limits, prohibition on autonomous
        power-off of safety-critical devices), and the corresponding automated tests.

  \item \textbf{Rule engine and recommendation engine.} User-defined automation rules, the
        AI-backed recommendation generator, cost-savings estimation, and the acceptance
        test cases from the test plan.

  \item \textbf{UI and integration.} React frontend pages, the capability-driven device
        control widget, the user flow, and final integration verification across all 65
        test cases.
\end{enumerate}

The team acknowledges that this was not a formally documented incremental plan with written
increment specifications. The division above reflects the actual sequence of verified
milestones rather than a pre-approved plan document, which is consistent with the
agile-style nature of the process.

% ────────────────────────────────────────────────────────────────────────────────
\section{Agile Practices Applied}
\label{sec:agile}

The course distinguishes between Agile as a \textit{philosophy} (the 2001 Manifesto) and
Scrum as a specific \textit{framework}. SHEO embodied several Agile values and principles
but did not run formal Scrum ceremonies. Each practised value is documented below.

\subsubsection*{Working software over documentation}
The most consistently applied Agile value was the primacy of a runnable, tested product.
At no point during development was the system reduced to a non-executing state; every
increment was required to leave the full test suite passing before integration. The 54
automated test cases serve as the continuous evidence of this discipline. Documentation
(the SDD, test plan, configuration management plan) was produced to capture design
decisions, but it was never treated as a substitute for a working system.

\subsubsection*{Short increments and continuous delivery of value}
Features were integrated in small, coherent units---a single endpoint with its test, a new
capability handler, a frontend page---rather than in large, infrequent releases. This kept
the feedback loop short and avoided the integration risk that accumulates when parallel
development streams diverge over extended periods.

\subsubsection*{Continuous testing}
Testing was not deferred to a distinct phase. Unit and integration tests were written
alongside the corresponding implementation, and the full suite was re-executed after every
non-trivial change. This is a direct application of the Agile principle that technical
excellence and good design enhance agility.

\subsubsection*{Responding to change}
Several design decisions were revised in response to feedback during development: the
original role model was renamed to align with the course's stakeholder taxonomy;
safety-rule NFRs were tightened after an adversarial review identified a cross-unit
access vulnerability; and the capability schema was generalised to accommodate additional
device types without modifying existing handlers. Each change was absorbed without
disrupting the passing test suite, demonstrating the architectural support for change that
Agile values require.

\subsubsection*{Definition-of-Done analogue}
Although no formal Scrum Definition of Done was documented, the team operated with an
implicit and consistently applied criterion: \textit{a feature is complete when (a)~it has
at least one passing automated test that exercises its primary execution path, and
(b)~the seeded demonstration database reflects its behaviour such that an evaluator can
observe it without additional setup.} This criterion operationalises two Scrum concepts---
the Increment and the Definition of Done---in a lightweight form appropriate to a small
student project.

\subsubsection*{Honest scope statement}
Formal Scrum roles (Product Owner, Scrum Master), artefacts (Product Backlog, Sprint
Backlog), and events (Sprint Planning, Daily Standup, Sprint Review, Retrospective) were
not adopted. The team was too small to benefit from the full Scrum overhead, and the
academic calendar imposed a fixed delivery date rather than an open-ended product roadmap.
The process is therefore characterised as \textit{agile-style incremental} rather than
Scrum.

% ────────────────────────────────────────────────────────────────────────────────
\section{Quality Throughout the Process}
\label{sec:quality}

The course cites ISO/IEC~9126~\cite{iso9126} as the quality standard for software products and emphasises
that quality assurance must run \textit{throughout} the development process, not only at
the end. Table~\ref{tab:iso9126} maps each ISO/IEC~9126 quality characteristic to the
concrete mechanism or non-functional requirement (NFR) that addressed it in SHEO, together
with the earliest phase at which that mechanism was active.

\begin{xltabular}{\linewidth}{@{} l Y Y l @{}}
\caption{ISO/IEC~9126 quality attributes mapped to SHEO mechanisms}\label{tab:iso9126}\\
\toprule
\textbf{Quality attribute} & \textbf{Definition (ISO/IEC~9126)} &
\textbf{SHEO mechanism / NFR} & \textbf{Active from} \\
\midrule
\endfirsthead
\toprule
\textbf{Quality attribute} & \textbf{Definition (ISO/IEC~9126)} &
\textbf{SHEO mechanism / NFR} & \textbf{Active from} \\
\midrule
\endhead
\textbf{Functionality} &
  The set of features present and their suitability for intended tasks. &
  Functional requirements REQ-4.1 through REQ-4.5 (monitoring, control, rules,
  recommendations, savings); validated by 65 automated acceptance and integration tests;
  traced in the Requirements Traceability Matrix. &
  Requirements \\
\addlinespace
\textbf{Reliability} &
  Ability to maintain performance under stated conditions over time; fault tolerance. &
  NFR-SAF: rate limits on device commands, prohibition on autonomous power-off of
  safety-critical appliances, input validation on all public endpoints. An adversarial
  review identified and closed a cross-unit data-isolation defect prior to release. &
  Design \\
\addlinespace
\textbf{Usability} &
  Ease of learning and use; appropriateness of the interface for target users. &
  Capability-driven UI widget: device controls are rendered from the capability schema,
  eliminating hard-coded per-device pages and ensuring consistent constraint enforcement
  (e.g.\ AC target temperature bounded to $16 \le T \le 30$\,\textdegree C). The seven
  UI golden rules were applied throughout UI design; a usability heuristic walkthrough
  was conducted before the frontend was built. &
  Design \\
\addlinespace
\textbf{Efficiency} &
  Resource consumption relative to performance under stated conditions. &
  NFR-PER: API response time target $\le 200$\,ms under simulated concurrent load;
  in-process database avoids network round-trips in the demo; the background simulator
  runs as a non-blocking asynchronous task to avoid blocking the request loop. &
  Design \\
\addlinespace
\textbf{Maintainability} &
  Ease of modifying the system to correct faults, improve performance, or adapt to a
  changed environment. &
  Capability schema as the single extension point for new device types (Open--Closed
  Principle); Strategy pattern in the adapter layer isolates device-specific logic;
  Repository pattern decouples persistence from business logic; all public interfaces
  documented with type hints and docstrings. &
  Design \\
\addlinespace
\textbf{Portability} &
  Ability to transfer the system to different environments. &
  Standard Python packaging, environment-variable configuration, and a
  containerisation-ready architecture (no hard-coded paths or OS-specific system calls);
  hardware abstracted behind the simulator interface so real IoT drivers can be substituted
  without changing application logic. &
  Design \\
\bottomrule
\end{xltabular}

\bigskip

The critical point is that none of these mechanisms was applied solely at the end of the
project. Reliability constraints were encoded as safety rules before the first device
endpoint was written. Maintainability was designed in from the outset through the capability
schema and layered architecture. Efficiency targets were stated as NFRs and informed
technology choices. Usability was evaluated iteratively through wireframes and heuristic
review before the frontend was built. This continuous, cross-phase attention to quality is
what both the course and the Agile philosophy prescribe, and it is what enabled SHEO to
reach a state of 65 passing tests, a fully seeded and demonstrable product, and no known
open defects at the point of submission.
